% Umbrella for the features; each feature is a \subsection in 06_ to 09_.
% Shared rules live here: unit of measurement (written), missing values (TODO).
% TODO statistics section: combine the per-artery predictions into one risk per participant, e.g.
%   arteries of one participant modelled as correlated, then P(any artery) = 1 - prod(1 - p_artery).
\section{Geometric Feature Extraction}
\label{sec:method-features}

Four geometric features were computed as candidate predictors of plaque development (Table~\ref{tab:method-features}).
Each was chosen for a published link between the geometry it measures and low wall shear stress, and all were computed from the geometry alone, since simulating flow through every tree of a population cohort is not feasible.

Three of the features were read from the labeled centerline tree of Section~\ref{sec:method-labeling}, in which every point carries a position in millimeters, a lumen radius and an artery label.
The first point of each tree is the ostium, where the artery leaves the aorta.
Every artery was traversed from it, so that distance along an artery and the generation of a bifurcation are both counted from the ostium and refer to the same anatomical starting point in every participant.
The fourth feature, lumen volume to myocardial mass, was computed from the segmentations of Section~\ref{sec:method-segmentation} rather than from the tree.

Each feature was computed per artery or per bifurcation rather than per participant, since wall shear stress acts locally and a single value per participant would average a harmful geometry in one artery with normal geometry in the others.
Curvature was computed for each of the LAD, the LCx and the RCA, and the bifurcation angle and daughter ratio for each bifurcation.
Lumen volume to myocardial mass is the exception, as it relates the lumen of the whole tree to the myocardium it supplies and therefore takes one value per participant.

\begin{table}[htbp]
  \centering
  \caption{Three of the four features are read from the labeled centerline tree; only lumen volume to myocardial mass needs the myocardium segmentation.}
  \label{tab:method-features}
  \begin{tabular}{llll}
    \toprule
    Feature & Measured on & Input & Section \\
    \midrule
    Turning-angle curvature & each named artery & centerline points & \ref{sec:method-curvature} \\
    Bifurcation angle & each bifurcation & daughter directions & \ref{sec:method-angle} \\
    Daughter ratio & each bifurcation & daughter radii & \ref{sec:method-ratio} \\
    Lumen volume to myocardial mass & whole tree & lumen and myocardium masks & \ref{sec:method-vm} \\
    \bottomrule
  \end{tabular}
\end{table}
